\chapter*{Abstract}
\addcontentsline{toc}{chapter}{Abstract}
% Written last, from the findings of Chapters 3 and 4. Every number is a macro.
After a stroke, one hand often stays weak. A robotic sixth finger worn on that hand can
help it grip, but it must move when the wearer wants it to. This project asked whether
brain signals recorded by an EEG cap can give that command when the patient imagines
moving the finger, and how much practice a new patient must provide before a model can
read them. We used a public recording of \nSubjects{} stroke patients who imagined moving
either the affected hand or the sixth finger.

An integrity audit first found recordings that appear twice in the release, and these
were excluded. The dataset paper reports about \pct{\rqOnePublishedGroupBCspLda} for
telling hand from sixth finger, with a script whose first step had already seen the test
trials. Our rebuild of the authors' recipe scored \pct{\rqOneCspLdaReplicaPaperReplicaMean}
that way and \pct{\rqOneCspLdaReplicaCvLeakfreeMean} once every step learned from the
training trials only, where guessing scores 50\,\%; a standard method fell from
\pct{\rqOneCspLdaPaperReplicaMean} to \pct{\rqOneCspLdaCvLeakfreeMean}. The same mistake
inflated two public datasets of healthy people as well.

In a calibration study planned and frozen in advance, one transfer configuration, in
which each of the new patient's windows counted as much as one of the other patients'
thousands, reached \pct{\rqTwoRestAhSfRctTslrKOneMean} on three choices (rest, hand,
sixth finger) after one practice trial of each movement, about \rqTwoMinutesKOne{}
minutes, against \pct{\rqTwoRestAhSfTsLrKOneMean} for the same model trained from scratch
on those four windows, which needed ten trials of each, about \rqTwoMinutesKTen{} minutes,
to catch up. Transfer then stayed nearly flat, and with the whole practice set the model
trained from scratch ended \rqTwoRestAhSfTsLrMinusRctTslrKAllDiff{} points higher, a
difference that was not tested
(\cipts{\rqTwoRestAhSfTsLrMinusRctTslrKAllDiffCiLo}{\rqTwoRestAhSfTsLrMinusRctTslrKAllDiffCiHi}).
For a new command, the confirmatory transfer method showed no clear benefit before therapy (it did after therapy,
so the result did not replicate), and a model trained on hand imagery alone detected
sixth-finger imagery at least as well, so no distinct sixth-finger command was shown. In
exploratory follow-ups, a switch for the sixth finger against rest was right
\pct{\exFourTsLrCentralMean} of the time window by window, but it also responded to
\pct{\exFourTsLrCentralFiresAhMean} of imagined hand movements, and played continuously it
caught only \pct{\exFiveFingerFingerOnMean} of the sixth-finger periods, with
\exFiveFingerFalsePerMinMean{} false switch-ons per minute of rest. Everything was tested
offline on recordings.

Finally, we built a demonstrator and website in which the trained model reads a real
patient's held-out recording in the browser and moves a drawn hand with a robotic sixth
finger. It runs at \url{\siteURL}.

\vspace{0.5em}
\noindent\textbf{Keywords:} brain--computer interface, EEG, motor imagery, stroke,
supernumerary finger, calibration, transfer learning, data leakage.
